\chapter{Online Learning: The Biggest Lever}
\label{ch:online}

If this book had to be compressed to one operational sentence, it would be:
\emph{in a drifting environment with delayed labels, a small daily gradient
step is worth more than any architecture you will try.} This chapter
quantifies that claim, derives the protocol, and maps the cliffs.

\section{The setting: delayed labels are an invitation}

The competition's lag mechanism --- yesterday's full responder table
delivered every morning (Chapter~\ref{ch:problem}) --- is not just a
feature source. It is a supervision stream: every day, 968 labeled steps
per symbol arrive, drawn from \emph{the most recent regime in existence}.
A frozen model ignores them; an adaptive one converts them into score.

The measured conversion rate, on our standard harness:

\begin{center}
\begin{tabular}{@{}lrrr@{}}
\toprule
model & static & online & $\Delta$ \\
\midrule
gru (single branch) & $+0.0025$ & $+0.0075$ & $+0.0050$ \\
gru\_modelr & $+0.0076$ & $+0.0089$ & $+0.0013$ \\
lstm\_modelr & $+0.0064$ & $+0.0092$ & $+0.0028$ \\
transformer & $+0.0047$ & $+0.0050$ & $+0.0003$ \\
sig-transformer (right lr) & $+0.0059$ & $+0.0062$ & $+0.0004$ \\
sig-transformer (RNN's lr) & $+0.0059$ & $+0.0044$ & $\bm{-0.0015}$ \\
\bottomrule
\end{tabular}
\end{center}

Three facts jump out, and the rest of the chapter unpacks them: the effect
is \textbf{large} (up to $+0.005$ --- compare the whole competitive range
of $0.012$); it is \textbf{family-dependent} (RNNs feast, attention
nibbles); and it can go \textbf{catastrophically negative} under
mis-tuning.

\section{The protocol}

One correct loop, matching deployment exactly
(Chapter~\ref{ch:validation}):

\begin{enumerate}
\item For each day $d$ of the evaluation walk, in calendar order:
\item \quad receive day $d{-}1$'s labels; form the day-batch
  (all symbols, all 968 steps);
\item \quad take one gradient step on the weighted-$R^2$ loss of that
  batch at learning rate $\eta_{\text{refit}}$, gradient-clipped;
\item \quad predict day $d$ causally; score; continue.
\end{enumerate}

Design choices inside that loop that we tested or inherited:

\begin{itemize}
\item \textbf{One step, not convergence.} The daily batch is one noisy
  draw from a moving distribution; fitting it well is overfitting it.
  A single clipped step is a \emph{tracking} update --- exponential
  smoothing in weight space.
\item \textbf{Which loss:} refit on the \emph{target} head only
  (the reference solution refits on \resp{6} alone, not the auxiliary
  heads). The aux heads did their job shaping the representation offline;
  online, you chase the target.
\item \textbf{Fresh optimizer state.} Our \code{update()} constructs the
  optimizer anew each day --- no stale Adam moments carried across a
  regime. With one step per day, momentum has nothing useful to
  accumulate anyway; a clean clipped step at the right $\eta$ is easier to
  reason about and measured fine.
\item \textbf{The scaler adapts too} (Chapter~\ref{ch:preprocessing}) ---
  but only from already-predicted days, and its setting is frozen during
  lr sweeps to keep the experiment single-variable.
\end{itemize}

\section{The learning-rate cliffs: a family-specific law}

The central empirical finding. Sweeping $\eta_{\text{refit}}$ per family:

\begin{center}
\begin{tabular}{@{}lrr@{}}
\toprule
$\eta_{\text{refit}}$ & lstm\_modelr & sig-transformer \\
\midrule
$0$ (static) & $+0.0063$ & $+0.0059$ \\
$10^{-5}$ & --- & $+0.0061$ \\
$3\times10^{-5}$ & --- & $\bm{+0.0062}$ \\
$10^{-4}$ & $+0.0083$ & $+0.0061$ \\
$3\times10^{-4}$ & $+0.0091$ & $+0.0044$ \\
$10^{-3}$ & $\bm{+0.0092}$ & $-0.0090$ \\
$3\times10^{-3}$ & $+0.0070$ & --- \\
\bottomrule
\end{tabular}
\end{center}

Read the columns separately, then together:

\begin{itemize}
\item Each family has a \textbf{plateau} (roughly a decade wide) flanked
  by cliffs: too small wastes the labels, too large lets one day's noise
  scramble the network ($-0.009$ is not a typo; that is a destroyed
  model).
\item The plateaus are \textbf{30$\times$ apart} between families
  ($10^{-3}$ vs $3\times10^{-5}$). A default tuned on one family sits on
  the other's cliff.
\item Consequence for research process: \textbf{no online-behavior verdict
  about an architecture is valid until its own lr sweep has run.} We
  spent two runs believing ``signature-transformers regress under refit''
  --- an architectural-sounding conclusion that was actually a borrowed
  hyperparameter (the full post-mortem is \S\ref{ch:signatures}).
\end{itemize}

\begin{warnbox}[: the borrowed-hyperparameter fallacy]
Whenever a new family underperforms under a procedure tuned on an old
family, the procedure's hyperparameters are confounded with the
architecture. Sweep before concluding. This fallacy has a long history in
deep learning (``X doesn't work'' $\to$ ``X needed its own lr'') and at
low SNR it is nearly undetectable without a deliberate sweep, because all
the numbers involved are small.
\end{warnbox}

\section{Why recurrence adapts best: a working hypothesis}

The pattern --- RNNs gain 10$\times$ more from identical updates --- wants
an explanation. Ours, offered as hypothesis with its supporting
observations: a gated RNN's function is \emph{smoothly} parameterized
(small weight changes $\to$ small, well-conditioned changes in its
temporal filtering), so a single clipped step moves it a small, useful
distance. Attention's computation routes through softmax comparisons;
identical step sizes can reorganize retrieval patterns discontinuously,
so safe steps must be tiny --- and tiny steps track drift slowly.
Consistent with this: the attention families' plateaus are both lower
\emph{and} narrower, and their best online gains ($+0.0004$) resemble a
strongly damped version of the RNN effect. Whatever the mechanism's true
name, the operational rule stands: \textbf{if online adaptation is
available, weight your architecture portfolio toward families that can
use it.}

\section{Variations we tested or queued}

\begin{itemize}
\item \textbf{Recency window + refit} (tested): the two compose --- train
  on the recent 280--400 days, \emph{and} refit daily. The window sweep's
  knee (50 $\to$ 400 days: $+0.0032 \to +0.0071$ online for the plain GRU)
  shifts \emph{up} at every length once refit is on; neither replaces the
  other.
\item \textbf{Refit-lr decay schedules} (queued): a per-day decay
  ($\eta_d = \eta_0 \cdot 0.5^{d/h}$) interpolates between adaptation and
  stability across a long test period; expected small positive by
  analogy, unproven here.
\item \textbf{Curriculum warm-starts across history} (tested, negative):
  chronologically warm-starting through 800 days of history scored
  \emph{below} the plain recent-window model ($+0.0077$ vs $+0.0089$).
  Old regimes are not a curriculum; they are interference
  (Chapter~\ref{ch:negative}).
\item \textbf{Nightly full retrains} (industry practice, out of
  competition scope): the heavyweight sibling of the daily step ---
  retrain on the trailing window every $k$ days. Composes with, does not
  replace, the light daily step.
\end{itemize}

\begin{jsbox}[: the lever in context]
Total spread from architecture choice among our serious single models:
$\sim 0.002$. Total spread from online refit on one architecture:
$\bm{0.003}$ (static $\to$ tuned online). Total spread from the refit
lr within one architecture: $\bm{0.018}$ ($-0.009$ to $+0.0092$).
The hyperparameter of the \emph{procedure} dwarfs the choice of the
\emph{model}. Allocate tuning attention accordingly.
\end{jsbox}

\begin{drillbox}
Extend Chapter~\ref{ch:rnn}'s drill: add a second architecture (any small
transformer) to your miniature drifting-panel setup and sweep refit lr
for both on a log grid from $10^{-5}$ to $10^{-2}$. Plot both curves on
one axis. You should see displaced plateaus and at least one
crossed-cliff configuration --- the borrowed-hyperparameter fallacy,
reproduced on your desk for the cost of twenty minutes of compute.
\end{drillbox}
